\documentclass[11pt]{article}

% --- Idioma y fuentes ------------------------------------------------
\usepackage[spanish,es-noshorthands,es-tabla]{babel}
\usepackage[T1]{fontenc}
\usepackage{lmodern}

% --- Geometría y encabezado ------------------------------------------
\usepackage[a4paper,margin=2.5cm]{geometry}
\usepackage{fancyhdr}
\setlength{\headheight}{38.22325pt}
\addtolength{\topmargin}{-6.22325pt}

% --- Matemáticas ------------------------------------------------------
\usepackage{amsmath,amssymb}

% --- Colores, tablas, listas, código ----------------------------------
\usepackage[dvipsnames,table]{xcolor}
\usepackage{booktabs}
\usepackage{array}
\usepackage{longtable}
\usepackage{enumitem}
\usepackage{listings}
\renewcommand{\arraystretch}{1.5}

\lstset{
  language=R,
  basicstyle=\ttfamily\footnotesize,
  commentstyle=\color{gray}\itshape,
  keywordstyle=\color{blue!70!black},
  stringstyle=\color{teal!60!black},
  breaklines=true,
  columns=fullflexible,
  showstringspaces=false
}


\pagestyle{fancy}
\fancyhf{}
\fancyhead[L]{\textbf{Econometría II}\\ Trabajo práctico --- Etapa 1}
\fancyhead[C]{Universidad de Santiago de Chile}
\fancyfoot[C]{\thepage}

\begin{document}
\sloppy

% =====================================================================
%  PORTADA
% =====================================================================
\begin{titlepage}
\centering
\vspace*{3cm}
{\large Universidad de Santiago de Chile}\\[6pt]
{\large Econometría II}\\[2cm]
{\LARGE \textbf{Etapa 1: Construcción de la base de datos}}\\[10pt]
{\large Trabajo Práctico --- Admisión Universidad 2026}\\[3cm]

\textbf{Integrantes:}\quad
David Cuevas $\bullet$ Javier Li-Yau $\bullet$ Benjamín Lizana $\bullet$ Ana Nacif

\vspace{0.6cm}

\textbf{Profesor:}\quad Brian Pfeng Ibacache
\qquad
\textbf{Fecha:}\quad 8 de octubre de 2026
\end{titlepage}

% =====================================================================
\subsection*{1 \; Introducción}
% =====================================================================

Con cuatro archivos del DEMRE del proceso de admisión 2026 armamos una
base de 108.251 estudiantes que egresaron de enseñanza media en 2025
como alumnos regulares y que postularon por la vía regular. De ellos,
60.300 (55,7\,\%) se matricularon por vía regular en una de sus tres
primeras preferencias. Los otros 47.951 (44,3\,\%) se matricularon en la
cuarta preferencia o más abajo (8.403) o no se matricularon (39.548). Nuestro trabajo apunta a explicar
de qué depende matricularse dentro de las tres primeras preferencias.

La base tiene una fila por estudiante y 30 variables. De cada uno trae
su promedio de notas, sus puntajes NEM y ranking, sus puntajes de
Competencia Lectora y Matemática 1, su sexo, su tramo de ingreso
familiar, si fue preseleccionado PACE o BEA y la preferencia en que se
matriculó. De su colegio de egreso trae la dependencia, la rama
educacional, la región y la comuna, y tres variables que calculamos
nosotros: inscritos en la PAES, matriculados y mediana del tramo de
ingreso de sus inscritos. Con esta información se estimará el modelo
logit de la Etapa 3.

\medskip\hrule\medskip

% =====================================================================
\subsection*{2 \; Fuentes de datos}
% =====================================================================

Todos los datos vienen del portal de bases de datos del DEMRE, proceso de
admisión 2026. Usamos cuatro archivos, que se unen por el identificador
del estudiante, \texttt{ID\_aux}.

\begin{center}
\begin{tabular}{>{\raggedright\arraybackslash}p{3.4cm}
                >{\raggedright\arraybackslash}p{6.4cm}
                >{\raggedleft\arraybackslash}p{2.6cm}}
\toprule
\textbf{Archivo} & \textbf{Contenido} & \textbf{Filas} \\
\midrule
B (\texttt{df1}) & Inscritos: colegio, domicilio, sexo, ingreso familiar & 317.700 \\
C (\texttt{df2}) & Notas, NEM, ranking y puntajes PAES & 317.700 \\
D (\texttt{df3}) & Postulaciones (una fila por carrera postulada) & 1.686.689 \\
Matrícula (\texttt{df4}) & Vía, tipo de matrícula, preferencia y puntaje & 128.345 \\
\bottomrule
\end{tabular}
\end{center}


\medskip\hrule\medskip

% =====================================================================
\subsection*{3 \; Construcción de la base}
% =====================================================================

\subsubsection*{3.1 \; Variables del colegio}

Antes de filtrar, calculamos las variables del colegio agrupando por
establecimiento (\texttt{RBD}). Se calculan antes porque necesitan a todos
los estudiantes del colegio, no solo a los que quedan en la base final.

\begin{itemize}[leftmargin=1.5em, itemsep=0.3em]
  \item Del archivo B: número de inscritos del colegio y mediana del
        tramo de ingreso per cápita familiar de sus inscritos.
  \item Del archivo de matrícula: número de matriculados por vía regular y
        tipo de matrícula 1 que egresaron del colegio. Los colegios sin
        matriculados quedan con 0.
\end{itemize}

Antes de agrupar hicimos dos ajustes. Primero, el código 99 de
\texttt{INGRESO\_PERCAPITA\_GRUPO\_FA} («prefiero no responder») se
convierte en dato faltante (\texttt{NA}) apenas se lee el archivo B, para
que no entre en el cálculo del colegio. Segundo, los inscritos sin
\texttt{RBD} no se agrupan: si se agruparan, R los trataría como un solo
colegio que no existe.

\subsubsection*{3.2 \; Selección de columnas y unión}

De cada archivo dejamos solo las columnas que usamos (detalle en la
sección 7). El archivo D trae una fila por carrera postulada; nos quedamos con las
postulaciones de tipo regular y lo dejamos en una fila por estudiante
(188.538 estudiantes). El archivo de
matrícula ya trae una fila por estudiante.

Luego unimos los cuatro archivos por \texttt{ID\_aux} con
\texttt{full\_join}, lo que da 317.700 filas, y a esa base le aplicamos
los filtros.

\subsubsection*{3.3 \; Filtros aplicados}

\begin{enumerate}[leftmargin=1.5em, itemsep=0.3em, label=\textbf{\alph*)}]
  \item \textbf{Promoción del año:} estudiantes que egresaron de enseñanza
        media en 2025 (\texttt{ANYO\_EGRESO} igual a 2025).
  \item \textbf{Alumno regular:} egresados de la promoción del año como
        alumnos regulares de un colegio chileno
        (\texttt{SITUACION\_EGRESO} igual a 1). Quedan fuera quienes
        validaron estudios o vienen del extranjero.
  \item \textbf{Postulación:} postularon al menos a una carrera por la
        vía regular (tienen \texttt{TIPO\_PREF} igual a REGULAR en el
        archivo D).
  \item \textbf{Matrícula:} se mantienen los que no se matricularon y,
        entre los matriculados, solo los de vía regular (\texttt{VIA}
        igual a 1) con tipo de matrícula «postulación regular»
        (\texttt{TIPO\_MATRICULA} igual a 1). Quedan fuera los
        matriculados por otra vía o con otro tipo de matrícula.
  \item \textbf{Inscripción en el proceso:} inscribieron al menos una
        prueba del proceso 2026 (\texttt{RINDIO\_PROCESO\_ACTUAL} igual
        a 1).
\end{enumerate}

No filtramos por preferencia: la preferencia en que se matriculó cada
estudiante es la que define la variable de respuesta (sección 4).

Además, sacamos a quienes no tienen puntaje NEM o lo tienen igual a 0,
porque sin notas de enseñanza media no se pueden comparar con el resto.

\subsubsection*{3.4 \; Observaciones en cada paso}

\begin{center}
\begin{tabular}{>{\raggedright\arraybackslash}p{9cm}
                >{\raggedleft\arraybackslash}p{3cm}}
\toprule
\textbf{Paso} & \textbf{Observaciones} \\
\midrule
Unión de los cuatro archivos & 317.700 \\
Egresados en 2025 & 218.989 \\
Alumnos regulares de la promoción del año & 218.066 \\
Postularon por la vía regular & 114.865 \\
No matriculados, o matriculados por vía regular y tipo 1 & 108.598 \\
Inscribieron alguna prueba del proceso 2026 & 108.596 \\
\rowcolor{yellow!30}
\textbf{Con puntaje NEM distinto de 0 (base final)} & \textbf{108.251} \\
\bottomrule
\end{tabular}
\end{center}

El filtro de puntaje NEM saca a 345 estudiantes.

\medskip\hrule\medskip

% =====================================================================
\subsection*{4 \; Variable de respuesta}
% =====================================================================

La variable de respuesta separa a quienes se matricularon en una de sus
tres primeras preferencias de quienes quedaron más abajo o no se
matricularon:

\[
\texttt{PREF\_TOP3\_DUMMY}_i =
\begin{cases}
1 & \text{si se matriculó en su primera, segunda o tercera preferencia} \\[4pt]
0 & \text{si quedó en su cuarta preferencia o más abajo, o no se matriculó}
\end{cases}
\]

\begin{center}
\begin{tabular}{lrr}
\toprule
\textbf{Valor} & \textbf{Estudiantes} & \textbf{Porcentaje} \\
\midrule
1 (preferencia 1 a 3)                          & 60.300 & 55,7\% \\
0 (preferencia 4 o más abajo, o sin matrícula) & 47.951 & 44,3\% \\
\midrule
\textbf{Total} & \textbf{108.251} & 100\% \\
\bottomrule
\end{tabular}
\end{center}

De los 60.300 con valor 1, 40.002 quedaron en la primera preferencia,
13.270 en la segunda y 7.028 en la tercera. De los 47.951 con valor 0,
8.403 se matricularon en la cuarta preferencia o más abajo y 39.548 no
se matricularon. Para estos últimos, \texttt{PREFERENCIA}, \texttt{VIA},
\texttt{TIPO\_MATRICULA} y \texttt{PTJE\_POND} quedan sin valor. Al
incluir a los no matriculados, la variable queda más equilibrada: 55,7\,\%
frente a 44,3\,\%.

\medskip\hrule\medskip

% =====================================================================
\subsection*{5 \; Variables creadas}
% =====================================================================

\subsubsection*{5.1 \; Variables del establecimiento}

Se pegan a cada estudiante por el \texttt{RBD} de su colegio de egreso.
Así cada estudiante tiene información de su colegio en el mismo proceso
2026.

\begin{itemize}[leftmargin=1.5em, itemsep=0.3em]
  \item \texttt{N\_ALUMNOS}: número de inscritos en el proceso 2026 que
        egresaron del colegio.
  \item \texttt{INGRESO\_COLEGIO}: mediana de
        \texttt{INGRESO\_PERCAPITA\_GRUPO\_FA} de los inscritos del
        colegio, sin contar a quienes respondieron 99. Usamos la mediana
        y no el promedio porque los tramos son una escala ordinal de
        ancho desigual: el último tramo es abierto hacia arriba, y
        promediar los números de los tramos supondría que la distancia
        entre ellos es siempre la misma. La variable queda en la escala de 1 a
        10; cuando un colegio tiene un número par de respuestas, la
        mediana puede terminar en ,5.
  \item \texttt{N\_MATRICULADOS}: número de egresados del colegio que se
        matricularon por vía regular y tipo de matrícula 1.
\end{itemize}

Los inscritos sin \texttt{RBD} quedan fuera de este cálculo (sección
3.1). En la base final no hay ninguno: casi todos los que no tienen
colegio registrado validaron estudios o vienen del extranjero, y salen con
el filtro de alumno regular. El colegio más grande tiene 947 inscritos y
447 matriculados.

No incluimos si el colegio es rural ni el número de alumnos por curso,
porque las bases del DEMRE no traen esa información: para construirlas
se necesitan los registros de establecimientos del MINEDUC.

\subsubsection*{5.2 \; Variables dummy}

\begin{itemize}[leftmargin=1.5em, itemsep=0.3em]
  \item \texttt{PREF\_TOP3\_DUMMY}: variable de respuesta (sección 4).
  \item \texttt{TP\_DUMMY}: 1 si el colegio es técnico profesional
        (rama T1 a T5), 0 si no, y sin valor si no se conoce la rama.
        Hay 18.877 estudiantes de colegios técnico profesionales y 89.374
        del resto.
  \item \texttt{PACE} y \texttt{BEA}: en el archivo B vienen con texto
        para los preseleccionados y vacías para el resto. Las
        convertimos en dummies: 1 si fue preseleccionado y 0 si no. Hay
        7.790 preseleccionados PACE y 9.916 BEA.
\end{itemize}

\medskip\hrule\medskip

% =====================================================================
\subsection*{6 \; Datos sin información}
% =====================================================================

\begin{center}
\begin{tabular}{>{\raggedright\arraybackslash}p{5.4cm}
                >{\raggedleft\arraybackslash}p{1.8cm}
                >{\raggedright\arraybackslash}p{5.6cm}}
\toprule
\textbf{Variable} & \textbf{Casos} & \textbf{Qué hicimos} \\
\midrule
\texttt{INGRESO\_PERCAPITA\_GRUPO\_FA} igual a 99 & 29.199 & Es «prefiero no responder», no un tramo: lo convertimos en \texttt{NA}. \\
\texttt{PACE} vacía & 100.461 & No fue preseleccionado: lo codificamos como 0. \\
\texttt{BEA} vacía & 98.335 & No fue preseleccionado: lo codificamos como 0. \\
\texttt{RBD} y variables del colegio & 0 & Los estudiantes sin colegio salen con el filtro de alumno regular y no entran al cálculo de las variables del colegio. \\
\texttt{PREFERENCIA}, \texttt{VIA}, \texttt{TIPO\_MATRICULA}, \texttt{PTJE\_POND} & 39.548 & No se matriculó: quedan sin valor y su variable de respuesta es 0. \\
\bottomrule
\end{tabular}
\end{center}

Fuera de las variables de matrícula de los no matriculados, la única
variable con datos faltantes es \texttt{INGRESO\_PERCAPITA\_GRUPO\_FA},
con 29.199 casos (27,0\,\% de la base). Si se deja así, el logit descartaría a esos estudiantes; para
la Etapa 3 evaluaremos tratar «prefiero no responder» como una categoría
más del tramo de ingreso.

\medskip\hrule\medskip

% =====================================================================
\subsection*{7 \; Definición de las variables}
% =====================================================================

Fuente: B, C, D y M (matrícula) son los archivos del DEMRE; «creada»
indica que la construimos nosotros.

\begin{longtable}{>{\raggedright\arraybackslash}p{5.4cm}
                  >{\raggedright\arraybackslash}p{1.4cm}
                  >{\raggedright\arraybackslash}p{6.6cm}}
\toprule
\textbf{Variable} & \textbf{Fuente} & \textbf{Definición} \\
\midrule
\endfirsthead
\toprule
\textbf{Variable} & \textbf{Fuente} & \textbf{Definición} \\
\midrule
\endhead
\bottomrule
\endfoot
\multicolumn{3}{l}{\textit{Identificación y respuesta}} \\
\texttt{ID\_aux} & B, C, D, M & Identificador único del estudiante. \\
\texttt{PREF\_TOP3\_DUMMY} & Creada & Variable de respuesta: 1 si se matriculó en una de sus tres primeras preferencias, 0 si en la cuarta o más abajo, o si no se matriculó. \\
\midrule
\multicolumn{3}{l}{\textit{Estudiante}} \\
\texttt{FECHA\_NACIMIENTO} & B & Mes y año de nacimiento, en formato MMAAAA. \\
\texttt{SEXO} & B & 1 hombre, 2 mujer. \\
\texttt{CODIGO\_COMUNA\_D} & B & Comuna de domicilio del estudiante. \\
\texttt{INGRESO\_PERCAPITA\_GRUPO\_FA} & B & Tramo de ingreso per cápita del grupo familiar, de 1 a 10. El 99 («prefiero no responder») se convirtió en \texttt{NA}. \\
\texttt{PACE} & B & 1 si fue preseleccionado PACE por el MINEDUC, 0 si no. \\
\texttt{BEA} & B & 1 si fue preseleccionado a la Beca de Excelencia Académica, 0 si no. \\
\texttt{ANYO\_EGRESO} & B & Año de egreso de enseñanza media (2025 en toda la base). \\
\texttt{SITUACION\_EGRESO} & B & Situación de egreso (1, promoción del año, alumno regular, en toda la base). \\
\texttt{RINDIO\_PROCESO\_ACTUAL} & B & 1 si inscribió al menos una prueba del proceso 2026 (1 en toda la base). \\
\midrule
\multicolumn{3}{l}{\textit{Notas y puntajes}} \\
\texttt{PROMEDIO\_NOTAS} & C & Promedio de notas de enseñanza media. \\
\texttt{PTJE\_NEM} & C & Puntaje NEM (notas de enseñanza media). \\
\texttt{PTJE\_RANKING} & C & Puntaje ranking de notas. \\
\texttt{CLEC\_REG\_ACTUAL} & C & Puntaje PAES de Competencia Lectora, rendición regular 2026. \\
\texttt{MATE1\_REG\_ACTUAL} & C & Puntaje PAES de Competencia Matemática 1, rendición regular 2026. \\
\midrule
\multicolumn{3}{l}{\textit{Postulación y matrícula}} \\
\texttt{TIPO\_PREF} & D & Tipo de postulación (REGULAR en toda la base). \\
\texttt{PREFERENCIA} & M & Preferencia en que se matriculó (1 a 20; sin valor si no se matriculó). Con ella se construye la variable de respuesta. \\
\texttt{VIA} & M & Vía de matrícula (1, regular, en todos los matriculados; sin valor si no se matriculó). \\
\texttt{TIPO\_MATRICULA} & M & Tipo de matrícula (1, postulación regular, en todos los matriculados; sin valor si no se matriculó). \\
\texttt{PTJE\_POND} & M & Puntaje ponderado con que se matriculó (sin valor si no se matriculó). \\
\midrule
\multicolumn{3}{l}{\textit{Colegio de egreso}} \\
\texttt{RBD} & B & Identificador del establecimiento. \\
\texttt{RAMA\_EDUCACIONAL} & B & H1 y H2 humanista científico (diurno y nocturno); T1 a T5 técnico profesional. \\
\texttt{GRUPO\_DEPENDENCIA} & B & 1 particular pagado, 2 particular subvencionado, 3 municipal, 4 servicio local de educación. \\
\texttt{CODIGO\_REGION} & B & Región del establecimiento. \\
\texttt{CODIGO\_COMUNA} & B & Comuna del establecimiento. \\
\texttt{TP\_DUMMY} & Creada & 1 si el colegio es técnico profesional (T1 a T5), 0 si no. \\
\texttt{N\_ALUMNOS} & Creada & Inscritos del proceso que egresaron del colegio. \\
\texttt{INGRESO\_COLEGIO} & Creada & Mediana de \texttt{INGRESO\_PERCAPITA\_GRUPO\_FA} de los inscritos del colegio, sin el código 99. \\
\texttt{N\_MATRICULADOS} & Creada & Egresados del colegio matriculados por vía regular y tipo de matrícula 1. \\
\end{longtable}

\texttt{PREFERENCIA} y \texttt{PTJE\_POND} quedan en la base como
información de la matrícula, pero no se usarán como variables
explicativas: la variable de respuesta se construye con
\texttt{PREFERENCIA}, y \texttt{PTJE\_POND} es el puntaje ponderado de
la carrera en que el estudiante ya quedó matriculado, así que depende
del resultado.

\medskip\hrule\medskip

% =====================================================================
\subsection*{8 \; Base final}
% =====================================================================

La base final (\texttt{df\_final}) tiene \textbf{108.251 filas y 30
columnas}, con una fila por estudiante: cada \texttt{ID\_aux} aparece una
sola vez. Los estudiantes vienen de 3.196 colegios distintos.

\begin{center}
\begin{tabular}{lrrrr}
\toprule
\textbf{Variable} & \textbf{Mínimo} & \textbf{Mediana} & \textbf{Promedio} & \textbf{Máximo} \\
\midrule
\texttt{PROMEDIO\_NOTAS}   & 4,20  & 6,23  & 6,19  & 7,00 \\
\texttt{PTJE\_NEM}         & 163   & 779   & 767   & 1.000 \\
\texttt{PTJE\_RANKING}     & 163   & 793   & 781   & 1.000 \\
\texttt{CLEC\_REG\_ACTUAL} & 0     & 634   & 640   & 1.000 \\
\texttt{MATE1\_REG\_ACTUAL}& 0     & 645   & 668   & 1.000 \\
\texttt{PTJE\_POND}        & 354,6 & 710,7 & 713,6 & 998,4 \\
\texttt{N\_ALUMNOS}        & 8     & 129   & 164   & 947 \\
\texttt{INGRESO\_COLEGIO}  & 1,0   & 4,0   & 4,8   & 10,0 \\
\texttt{N\_MATRICULADOS}   & 0     & 54    & 70    & 447 \\
\bottomrule
\end{tabular}
\end{center}

\texttt{PTJE\_POND} se resume solo para los 68.703 matriculados. Los
mínimos de 0 en \texttt{CLEC\_REG\_ACTUAL} y \texttt{MATE1\_REG\_ACTUAL}
corresponden a quienes no rindieron esa prueba (77 y 161 estudiantes).
En la base hay 59.799 mujeres (55,2\%) y 48.452 hombres (44,8\%).

\medskip\hrule\medskip

% =====================================================================
\subsection*{9 \; Conclusión}
% =====================================================================

La construcción de la base de datos permitió integrar en un único data
frame la información académica, socioeconómica, de postulación, matrícula
y establecimiento educacional de los estudiantes participantes del
proceso de admisión 2026.

Luego de aplicar los filtros establecidos en el trabajo, la muestra final
quedó compuesta por 108.251 estudiantes egresados de enseñanza media en
2025 como alumnos regulares que postularon por la vía regular: 68.703 se
matricularon por vía regular y bajo postulación regular, y 39.548 no se
matricularon. Asimismo, se incorporaron variables construidas a
nivel de establecimiento y variables dummy que permitirán facilitar el
análisis econométrico posterior. Los problemas de la base se corrigieron
antes de construirlas: el código 99 del ingreso se convirtió en
\texttt{NA}, el ingreso del colegio se calcula con la mediana y los
estudiantes sin colegio registrado no forman un colegio ficticio.

La variable de respuesta definida corresponde a
\texttt{PREF\_TOP3\_DUMMY}, la cual identifica si el estudiante logró
matricularse en una de sus tres primeras preferencias o, en cambio,
quedó en la cuarta o más abajo o no se matriculó. En la base final, un
55,7\,\% de los estudiantes se matriculó dentro de sus tres primeras
preferencias, mientras que un 44,3\,\% quedó más abajo o no se
matriculó.

En términos generales, la base construida cumple con las condiciones
establecidas para esta primera etapa y entrega una estructura adecuada
para avanzar hacia el análisis descriptivo y la posterior estimación de
un modelo Logit, con el objetivo de estudiar qué características
individuales, académicas y del establecimiento influyen en la
probabilidad de matricularse dentro de las tres primeras preferencias.

% =====================================================================
\end{document}
